\subsection{前测度的傅里叶变换}

设 E 是局部凸空间，且 \(\mu = (\mu_{\mathrm{V}})_{\mathrm{V} \in \mathcal{F}(\mathrm{E})}\) 是 E 上的前测度。对 E 上的每个连续线性形式 \(x'\)，记 \(\mu_{x'}\) 为前测度 \(x'\) 经 \(\mu\) 作用后在 R 上得到的测度。称 \(\mu\) 的傅里叶变换为定义在 \(F\mu\) 上的函数 \(E'\)，由下式给出

\[
(\mathcal {F} \mu) (x ^ {\prime}) = \int_ {\mathbf {R}} e ^ {i t} d \mu_ {x ^ {\prime}} (t).\tag{2}
\]

设 \(\lambda\) 是 E 上的有界测度。称 \(\lambda\) 的傅里叶变换为定义在 \(\mathbf{E}'\) 上的函数，由下式给出

\[
(\mathcal {F} \lambda) (x ^ {\prime}) = \int_ {\mathrm{E}} e ^ {i \langle x, x ^ {\prime} \rangle} d \lambda (x).\tag{3}
\]

设 \(\mu\) 是与 \(\lambda\) 对应的前测度。对每个 \(x' \in E'\)，定义在 \(\mu_{x'}\) 上的测度 \(\mathbf{R}\) 是在映射 \(x': \mathrm{E} \to \mathbf{R}\) 下测度 \(\lambda\)（定义在 \(\mathrm{E}\) 上）的像；由公式 (2) 和 (3) 立即推出 \(\mathcal{F}\mu = \mathcal{F}\lambda\)。

设 \(\mu\) 是 \(\mathbf{E}\) 上的任意前测度，且 \(u\) 是从 \(\mathbf{E}\) 到局部凸空间 \(\mathbf{E}_1\) 的连续线性映射。记 \(^t u\) 为从 \(\mathbf{E}_1'\) 到 \(\mathbf{E}'\) 的线性映射，它是 \(u\) 的转置；并记 \(\nu\) 为前测度 \(u(\mu)\)，定义在 \(\mathbf{E}_1\) 上。对每个 \(x_1' \in \mathbf{E}_1'\)，有 \(^t u(x_1') = x_1' \circ u\)，从而

\[
\nu_ {x _ {1} ^ {\prime}} = x _ {1} ^ {\prime} (\nu) = x _ {1} ^ {\prime} (u (\mu)) = \left(x _ {1} ^ {\prime} \circ u\right) (\mu) = \mu^ {t} u \left(x _ {1} ^ {\prime}\right).
\]

因此，

\[
\mathscr {F} (u (\mu)) = (\mathscr {F} \mu) \circ {} ^ {t} u.\tag{4}
\]

特别地，对 \(p_{V}\) 取从 E 到 E/V 的典范映射 \(V \in \mathcal{F}(E)\)（\(p_{\mathrm{V}}(\mu)\)）。定义在 E/V 上的前测度 \(\mu_{V}\) 与测度 \(^{t}p_{V}\) 对应，而 \(V^{\circ}\) 将 E/V 的对偶同构地映到 E' 中与 V 正交的子空间 \(E'\)。若借助 \((\mathrm{E}/\mathrm{V})'\) 将 \(V^{\circ}\) 与 \(^{t}p_{V}\) 识别，则

\[
(\mathcal {F} \mu) (x ^ {\prime}) = \int_ {\mathrm{E} / \mathrm{V}} e ^ {i \langle x, x ^ {\prime} \rangle} d \mu_ {\mathrm{V}} (x)\tag{5}
\]

对 \(x' \in V^{\circ}\) 都有。\(E' = \bigcup_{V \in \mathcal{F}(E)} V^{\circ}\)，所以前述公式刻画了函数 \(\mathcal{F}\mu\) 在 \(E'\) 上。最后，在 (5) 中置 \(x' = 0\)，可见 \(\mu\) 的总质量等于 \((\mathcal{F}\mu)(0)\)。

命题 3. --- 设 E 是局部凸空间。映射 \(\mu \mapsto \mathcal{F}\mu\) 从 E 上的前测度集合到 E' 上的函数集合是单射。

公式 (5) 使问题归结为 E 有限维的情形；由于每个有限维空间都同构于某个 \(R^{n}\)，甚至不妨设存在整数 \(n \geqslant 0\)，使得 \(E = R^{n}\)。因此需要证明：若 \(\mu\) 是 \(R^{n}\) 上的有界测度（不一定为正），且

\[
\int_ {\mathbf {R} ^ {n}} e ^ {i \langle x, y \rangle} d \mu (x) = 0
\]

对每个线性形式 \(y\)（其定义域为 \(\mathbf{R}^n\)），上述等式都成立，则 \(\mu = 0\)。

对每个整数 \(m \geqslant 0\)，令 \(\mathbf{G}_m\) 为 \(m \cdot \mathbf{Z}^n\) 在 \(\mathbf{R}^n\) 中的子群。记 \(\mathcal{C}_m\) 为连续函数 \(f\) 所成的向量空间，其中 f 定义在 \(\mathbf{R}^n\) 上并满足 \(f(x + a) = f(x)\) 对 \(x \in \mathbf{R}^n\) 和 \(a \in \mathbf{G}_m\) 成立。由 GT, X, §4, No.~4 的命题 8，\(\mathcal{C}_m\) 中每个函数都是形如 \(x \mapsto e^{2\pi i\langle x,q\rangle}\)（其中 \(q \in m^{-1} \cdot \mathbf{Z}^n\)）的函数的有限线性组合的一致极限。因此，\(\mu(f) = 0\) 对每个函数 \(f \in \mathcal{C}_m\) 成立。

设 \(f\) 是 \(\mathbf{R}^n\) 上具有紧支集的连续函数。对每个整数 \(m \geqslant 0\)，置 \(f_m(x) = \sum_{q \in G_m} f(x + q)\)。显然，对每个 \(x \in \mathbf{R}^n\)，上述级数只有有限项，且 \(f_m\) 属于 \(\mathcal{C}_m\)。此外，容易看出序列 \((f_m)\) 在每个紧集上一致收敛到 \(f\)，并且存在常数 \(C \geqslant 0\)，使得 \(|f_m| \leqslant C\) 对所有 \(m\) 成立。因此，由 §5 第 6 号命题 12，\(\mu(f) = \lim_{m \to \infty} \mu(f_m)\)。由于 \(f_m \in \mathcal{C}_m\)，有 \(\mu(f_m) = 0\)，最终得到 \(\mu(f) = 0\)。故 \(\mu = 0\)。

*备注。 --- 当 E 有限维时，E 的每个特征都具有形式 \(x \mapsto e^{i\langle x,x' \rangle}\)，其中 \(x' \in \mathrm{E}'\)（Théor. spect., 第 II 章 §1, No.~9, 命题 12 的推论 3）。此时，命题 3 由傅里叶变换的唯一性定理得出（同上，§1, No.~6, 命题 6 的推论）。*

